\documentclass{article}

\usepackage{amsmath,amssymb}
\usepackage{xurl}
\usepackage[hidelinks]{hyperref}

% Shared commands for notes in docs/paper-gaps/.

\DeclareUnicodeCharacter{2080}{\ifmmode{}_{0}\else\textsubscript{0}\fi}
\DeclareUnicodeCharacter{2081}{\ifmmode{}_{1}\else\textsubscript{1}\fi}
\DeclareUnicodeCharacter{2082}{\ifmmode{}_{2}\else\textsubscript{2}\fi}
\DeclareUnicodeCharacter{2099}{\ifmmode{}_{n}\else\textsubscript{n}\fi}

\newcommand{\Lean}{\textsc{Lean}}
\ExplSyntaxOn
\NewDocumentCommand{\leanid}{m}{
  \ifmmode
    \textsf{\detokenize{#1}}
  \else
    \group_begin:
    \urlstyle{sf}
    \tl_set:Nn \l_tmpa_tl {#1}
    \tl_replace_all:Nnn \l_tmpa_tl {\_} {_}
    \exp_args:NV \path \l_tmpa_tl
    \group_end:
  \fi
}
\ExplSyntaxOff
\providecommand{\tr}{\operatorname{tr}}
\newcommand{\tnleanIssueBase}{https://github.com/LionSR/TNLean/issues/}
\newcommand{\tnleanPrBase}{https://github.com/LionSR/TNLean/pull/}
\newcommand{\tnleanDocBase}{https://sirui-lu.com/TNLean/docs/}
\newcommand{\ghissue}[1]{\href{\tnleanIssueBase#1}{GitHub issue \##1}}
\newcommand{\ghpr}[1]{\href{\tnleanPrBase#1}{PR \##1}}
\newcommand{\doclink}[1]{\href{\tnleanDocBase#1.html}{\path{#1}}}

% Dirac notation and the identity operator, as in blueprint/src/macros/common.tex.
% \providecommand keeps note-local definitions authoritative.
\usepackage{dsfont}
\providecommand{\ket}[1]{|#1\rangle}
\providecommand{\bra}[1]{\langle #1|}
\providecommand{\braket}[2]{\langle #1 | #2 \rangle}
\providecommand{\ketbra}[2]{|#1\rangle\!\langle #2|}
\providecommand{\Id}{\mathds{1}}
\providecommand{\one}{\mathds{1}}
\providecommand{\C}{\mathbb{C}}
\providecommand{\MN}[1]{M_{#1}(\C)}
\providecommand{\MD}{M_D(\C)}

% External referencing for paper sources.  The build helper writes sanitized
% auxiliary files under build/paper-aux/ and excludes duplicated source labels.
\usepackage{xr}

% Import a generated paper auxiliary file with a caller-supplied label prefix.
% The candidate paths support builds from docs/paper-gaps/, the repository root,
% and build/paper-aux/ itself.
\newcommand{\importpaperaux}[2]{%
  \IfFileExists{../../build/paper-aux/#2.aux}{%
    \externaldocument[#1]{../../build/paper-aux/#2}%
  }{%
    \IfFileExists{build/paper-aux/#2.aux}{%
      \externaldocument[#1]{build/paper-aux/#2}%
    }{%
      \IfFileExists{#2.aux}{%
        \externaldocument[#1]{#2}%
      }{}%
    }%
  }%
}

% General external reference macros
\newcommand{\paperref}[2]{\ref{#1-#2}}
\newcommand{\papereqref}[2]{(\ref{#1-#2})}

% CPSV16 source (1606.00608 / MPDO-22-12-17-2.tex) external document and shortcuts
\importpaperaux{cpsv16-}{1606.00608/MPDO-22-12-17-2-tnlean}
\newcommand{\cpsvref}[1]{\paperref{cpsv16}{#1}}
\newcommand{\cpsveqref}[1]{\papereqref{cpsv16}{#1}}

% Verdict marker: \gapnote{<kind>}{<status>}. Kind is the policy
% classification (clarification, local-correction, scope-restriction,
% unfaithful, false-source, open-gap); status is open, wip (elimination
% underway), resolved, or historical. Machine-read by the site index;
% prints a verdict line.
\newcommand{\gapnote}[2]{\par\noindent\textbf{Verdict:} #1 (#2).\par}

\importpaperaux{fbc25-}{2502.20257/main-tnlean}

\title{Defect Covariance in State-Level Gauging}
\author{}
\date{2026-09-03}

\begin{document}
\maketitle
\gapnote{scope-restriction}{open}

\section{Source statement}

The modified-fusion lemma of
Ref.~\cite[lines~4215--4326]{FrancoRubio2025MPUGauging} constructs unitary
fusion operators from locally orthogonal matrix-product-state blocks. The
resulting modified Gauss operators leave the gauged state invariant
\cite[lines~4327--4335]{FrancoRubio2025MPUGauging}. Thus the paper derives the
local defect support and covariance used in the invariance argument from its
concrete defect tensors and fusion maps.

\section{Current formal scope}

The present universal-vector theorem starts with defect subspaces
$\mathcal K_{a,b}$, prescribed maps $D_{a,b}$, and matter states
$\ket{\Psi_\alpha}$. It assumes that every two-site matter slice lies in the
corresponding defect subspace and that, for every window and $g\in G$,
\begin{align}
    (D_{a,b}\otimes\Id)\ket{\Psi_\alpha}
     & =(D_{ag^{-1},gb}\otimes\Id)
    \ket{\Psi_{T_{j,g}\alpha}}.
\end{align}
Under these hypotheses it proves invariance for every unitary completion of the
prescribed maps. This is a conditional abstraction of the paper's argument,
not yet a derivation of the hypotheses from locally orthogonal MPS blocks.

\section{Elimination plan}

To recover the full source construction, instantiate the defect subspaces,
prescribed maps, and matter states from the FBC25 defect tensors; prove local
support and the displayed covariance from local orthogonality and the modified
fusion lemma; then specialize the abstract universal-vector theorem. The
four-domain CZX instance and this source-to-abstraction bridge remain tracked
by \ghissue{7569}.

\bibliographystyle{alpha}
\bibliography{references}

\end{document}